\section{Actual weighted DVR charts of an original regular local surface}

Let $S$ be a Noetherian local domain with maximal ideal
$\mathfrak m=(\pi,z)$. Assume that $\pi$ is an original prime element and that
$z\notin(\pi)$. These are actual regular parameters; no valuation, normalized
chart, dominance, fraction-field comparison, or desired order formula is an
input. Fix an original integer $N>0$ and put $c=z^N$.

\begin{definition}[The same literal weighted surface chart]
\label{def:gauss-regular-surface-weighted-chart}
\lean{EnvelopeProperty.GaussRegularSurfaceValuations.relation,
EnvelopeProperty.GaussRegularSurfaceValuations.chart,
EnvelopeProperty.GaussRegularSurfaceValuations.coefficientMap,
EnvelopeProperty.GaussRegularSurfaceValuations.chartParameter}
\leanok
Use the actual original polynomial and quotient
\[
 F_N=\pi U-z^N\in S[U],\qquad B_N=S[U]/(F_N).
\]
The coefficient map $S\to B_N$ and the element $u_N\in B_N$ are the actual
quotient maps of the original constants and $U$. Their original relation is
$\pi u_N=z^N$.
\end{definition}
\begin{proof}
This is precisely the genuine polynomial quotient \texttt{AdjoinRoot} of
$F_N$, with its original coefficient map and original root. The quotient
annihilates $F_N$, giving the displayed literal coefficient identity.
No integral closure or replacement of the original surface ring is used.
\end{proof}

\begin{lemma}[Original prime-parameter saturation and actual chart domain]
\label{lem:gauss-regular-surface-chart-domain}
\uses{def:gauss-regular-surface-weighted-chart}
\lean{EnvelopeProperty.GaussRegularSurfaceValuations.relation_saturated,
EnvelopeProperty.GaussRegularSurfaceValuations.chart_isDomain}
\leanok
The actual ideal $(F_N)$ is saturated under multiplication by every original
power of $\pi$, and the SAME literal quotient $B_N$ is a domain. Evaluation
$U\mapsto z^N/\pi$ embeds it in the original $S[1/\pi]$, and then in the SAME
original fraction field of $S$. In particular the original coefficient map
is injective.
\end{lemma}
\begin{proof}
The original prime property and $z\notin(\pi)$ give $z^N\notin(\pi)$.
Suppose $\pi P=F_NQ$ in the actual polynomial ring. Reduce every original
coefficient modulo $(\pi)$. The original quotient $S/(\pi)$ is a domain and
the image of $z^N$ is nonzero. Thus $0=-\overline{z^N}\,\overline Q$ forces
$\overline Q=0$. Every original coefficient of $Q$ is therefore divisible by
$\pi$, so $Q=\pi Q_1$ in $S[U]$. Cancel the original nonzero $\pi$ in this
domain to obtain $P=F_NQ_1$. Iterate this implication for powers of $\pi$.

Localize the ORIGINAL coefficients away from $\pi$. There $F_N$ is the
original unit $\pi$ times the monic polynomial $U-z^N/\pi$. Its quotient is
therefore the same domain $S[1/\pi]$, by actual evaluation and division by a
monic linear polynomial. Polynomial localization commutes with this original
coefficient localization. If an original polynomial maps to zero under the
stated evaluation, its membership in the localized relation ideal clears
by a genuine power of $\pi$. The saturation just proved cancels that power,
so its original kernel is exactly $(F_N)$. This gives the actual asserted
embedding and hence the domain property of the SAME quotient $B_N$.
\end{proof}

\begin{lemma}[The actual closed-centre prime of the original chart]
\label{lem:gauss-regular-surface-chart-closed-prime}
\uses{def:gauss-regular-surface-weighted-chart,lem:gauss-regular-surface-chart-domain}
\lean{EnvelopeProperty.GaussRegularSurfaceValuations.closedPrime,
EnvelopeProperty.GaussRegularSurfaceValuations.chart_residueIso,
EnvelopeProperty.GaussRegularSurfaceValuations.closedPrime_isPrime}
\leanok
The original extended ideal $\mathfrak p_N=\mathfrak mB_N$ is prime, and its
actual quotient, with the original coefficient and variable maps, is
\[
 B_N/\mathfrak p_N\simeq\kappa(S)[U].
\]
In particular $u_N\notin\mathfrak p_N$, while the original images of $\pi,z$
belong to $\mathfrak p_N$.
\end{lemma}
\begin{proof}
Apply the genuine double polynomial quotient isomorphism to the original
ideal $\mathfrak m$ and the literal relation $F_N$. Both $\pi,z$ belong to
$\mathfrak m$, and $N>0$, so the image of $F_N$ in $\kappa(S)[U]$ is the zero
polynomial. Its additional quotient by zero is canonically itself. The
isomorphism retains EVERY original constant and sends the original $u_N$ to
$U$. This target is a domain, and $U$ is nonzero in it. Therefore the actual
extended ideal is prime and does not contain $u_N$. The two parameter
memberships are the original ideal-map memberships.
\end{proof}

\begin{theorem}[An internally constructed original DVR with exact parameter orders]
\label{thm:gauss-regular-surface-original-dvr}
\uses{lem:gauss-regular-surface-chart-domain,lem:gauss-regular-surface-chart-closed-prime}
\lean{EnvelopeProperty.GaussRegularSurfaceValuations.localRing,
EnvelopeProperty.GaussRegularSurfaceValuations.localCoefficientMap,
EnvelopeProperty.GaussRegularSurfaceValuations.localParameter,
EnvelopeProperty.GaussRegularSurfaceValuations.local_maximalIdeal,
EnvelopeProperty.GaussRegularSurfaceValuations.localRing_isDVR,
EnvelopeProperty.GaussRegularSurfaceValuations.localParameter_irreducible,
EnvelopeProperty.GaussRegularSurfaceValuations.parameter_addVal,
EnvelopeProperty.GaussRegularSurfaceValuations.vertical_addVal,
EnvelopeProperty.GaussRegularSurfaceValuations.localCenter,
EnvelopeProperty.GaussRegularSurfaceValuations.fractionEmbedding,
EnvelopeProperty.GaussRegularSurfaceValuations.fractionEmbedding_injective,
EnvelopeProperty.GaussRegularSurfaceValuations.fractionEmbedding_coefficient,
EnvelopeProperty.GaussRegularSurfaceValuations.localFractionRing}
\leanok
The SAME actual localization $V_N=(B_N)_{\mathfrak p_N}$ is a Noetherian local
domain and a discrete valuation ring. Its actual maximal ideal is generated
by the original image of $z$, which is a genuine uniformizer. With the native
normalized additive valuation of this constructed DVR,
\[
 \operatorname{ord}_{V_N}(z)=1,\qquad
 \operatorname{ord}_{V_N}(\pi)=N.
\]
Its original coefficient map is centred at the SAME $\mathfrak m$, and its
actual fraction field is the SAME original $\operatorname{Frac}S$, through
an injective actual map retaining every original coefficient and the relation
$u_N=z^N/\pi$. No DVR or valuation-family certificate is supplied.
\end{theorem}
\begin{proof}
The polynomial quotient $B_N$ is Noetherian and, by the preceding actual
kernel argument, a domain. Its prime localization $V_N$ inherits both
properties and is local. The original $u_N$ avoids $\mathfrak p_N$, so its
actual localization image is a unit. The actual maximal ideal of $V_N$ is
$\mathfrak p_NV_N=(\pi,z)V_N$. The literal relation $\pi u_N=z^N$ and $N>0$
show that $\pi$ belongs to the ideal generated by $z$. Thus the SAME maximal
ideal is $(z)$. The image of $z$ is nonzero by the original domain embeddings,
and belongs to the maximal ideal, so the local domain is not a field. The
Noetherian local principal-maximal-ideal criterion now derives that THIS
$V_N$ is a DVR, and its original $z$ is irreducible, hence a uniformizer.
Its true native order is one. Apply the native additive valuation to the
original relation: the actual unit $u_N$ has order zero, so the order of
$\pi$ equals the order of $z^N$, namely $N$.

Every original member of $\mathfrak m$ maps into $\mathfrak p_N$ and hence
into the actual localization maximal ideal. Every original element outside
$\mathfrak m$ is a unit in $S$ and remains a unit in $V_N$. This proves the
WHOLE exact centre equality. The embedding of $B_N$ into the original
$\operatorname{Frac}S$ extends to its actual prime localization, since every
original denominator there is nonzero and hence maps to a unit in that
field. Its injectivity follows from the actual embedding and the localization
kernel criterion. Finally every element of the original $\operatorname{Frac}S$
is a ratio of two original elements of $S$, hence of their actual images in
$V_N$. This establishes the actual fraction-ring structure with that SAME
map and SAME original coefficients.
\end{proof}

The present construction supplies actual DVRs and exact original parameter
orders. Pointwise approximation of the vertical valuation, normalized fixed
Gauss coefficient values, convergence and residual difference quotients are
subsequent statements; they are not assumed or claimed by this construction.
